\documentclass[multi=tikzpicture,border=9pt]{standalone}
\usepackage[UTF8]{ctex}
\usepackage{amsmath,amssymb}
\usepackage{tikz}
\usetikzlibrary{arrows.meta,calc,patterns,decorations.pathmorphing}
\definecolor{main}{HTML}{17365D}
\definecolor{accent}{HTML}{AC3931}
\tikzset{every picture/.style={font=\small,line width=.75pt},
>={Stealth[length=2mm]},arr/.style={->},
box/.style={draw=main,rounded corners=2pt,fill=main!4,text width=3.6cm,minimum height=1cm,align=center},
note/.style={align=center,text width=9cm}}
\begin{document}

\begin{tikzpicture}
\draw[fill=main!3] (-3.5,-.65)--(3.5,-.65)--(4.5,.7)--(-2.5,.7)--cycle;
\draw[main,thick,rotate around={-32:(0,2)}] (0,2) ellipse (3.774 and .65);
\draw[accent,very thick] (0,2)--(3.2,0);\draw[dashed] (0,2)--(0,0);\draw (0,0)--(3.2,0);
\fill (0,2) circle (3pt);\fill[accent] (3.2,0) circle (3pt);\fill (0,0) circle (2pt);\node[above left] at (0,2){$C$};\node[below] at (0,0){$O$};\node[right] at (3.2,0){$A$};
\draw[arr] (0,0)--(0,4.5) node[above]{$z$};\draw[arr] (0,0)--(4,0) node[right]{$x$};\draw[arr] (0,0)--(-1.4,.8) node[above]{$y$};\draw[arr,accent] (-2.5,-.2)--(-3.9,.6) node[above]{$\mathbf B$};
\node[note] at (.3,-1.9){初始接触方向示意；$CA$为半径，$CO=r\sin\theta$\\固定引线接 $C$ 与平面，图示未画外部仪表};
\end{tikzpicture}

\begin{tikzpicture}
\draw[main,thick] (0,0) circle (2.2);\fill (0,0) circle (2pt);\node[below] at (0,0){$O$};\draw[arr] (0,0)--(2.2,0) node[midway,below]{$OA$};\fill[accent] (2.2,0) circle (3pt);\node[right] at (2.3,0){$A$};\draw[arr,accent] (2.2,0)--(2.2,1.4) node[right]{$\dot{\mathbf A}_{\rm geom}$};\node[box] at (-3.5,1.3){接触位置在移动};\node[box] at (-3.5,-1.3){接触材料速度为零};\node[note] at (-.8,-3.4){无滑动更换接触材料点；不能把两种速度认作同一量};
\end{tikzpicture}

\begin{tikzpicture}
\draw[arr] (0,0)--(4.5,0) node[right]{$\hat e$};\draw[arr] (0,0)--(0,4.2) node[above]{$\hat z$};\fill (0,2.4) circle (3pt);\node[left] at (0,2.4){$C$};\fill[accent] (3.2,0) circle (3pt);\node[below] at (3.2,0){$A$};\draw[accent,very thick] (3.2,0)--(-.8,3);\draw[arr,main,thick] (0,2.4)--(-1.4,3.45) node[above]{$\boldsymbol\Omega$};\node[right] at (1.6,1.2){$\mathbf v_{CA}=0$};\node[note] at (1.2,-1.3){$\boldsymbol\Omega=-\omega\sin\theta\,\hat u$\\$|\boldsymbol\Omega|=|\omega|\sin\theta$};
\end{tikzpicture}

\begin{tikzpicture}
\node[box] (a) at (0,2){圆盘径线 $CA$\\$\mathbf v_{\rm mat}=0$};\node[box] (b) at (5,2){平面与引线\\$\mathbf v_{\rm mat}=0$};\node[box] (c) at (2.5,0){真实电动势\\$\mathcal E=0$};\draw[arr] (a)--(c);\draw[arr] (b)--(c);\node[box] at (2.5,-2){几何磁通导数\\一般非零};\node[note] at (2.5,-3.6){几何轮廓不断选取新的材料路径，不能无条件用磁通导数替代电动势};
\end{tikzpicture}

\begin{tikzpicture}
\node[box] (n0) at (0,0) {无滑动条件\\$\mathbf v_C=\mathbf v_A=0$};
\node[box] (n1) at (5,0) {瞬时零速度轴\\$\boldsymbol\Omega\parallel CA$};
\node[box] (n2) at (5,-2.3) {力积分定义\\$\mathcal E=\oint f_{\parallel}\,d\ell/q$};
\node[box] (n3) at (0,-2.3) {题设结论\\$\mathcal E=W=0$};
\draw[arr] (n0)--(n1);\draw[arr] (n1)--(n2);\draw[arr] (n2)--(n3);
\node[note] at (2.5,-4.1) {几何回路磁通变化不等于静态磁场中材料路径的动生电动势};
\end{tikzpicture}
\end{document}
